\documentclass[preview]{standalone}
\usepackage{amsmath}
\usepackage{amssymb}
\usepackage{stellar}
\usepackage{definitions}
\usepackage{bettelini}
\begin{document}
\id{diffeq-functional-setting}
\genpage
\section{Function spaces}

\includesnpt{bounded-continuous-functions-uniform-norm}

\begin{snippetdefinition}{diffeq-function-space-norms}{Uniform derivative norms}
    For \(A\subseteq\realnumbers\), let \(B(A,\realnumbers^m)\) be the space
    of bounded \function[functions], with
    \(\|u\|_{\infty,A}=\sup_{x\in A}\|u(x)\|\).
    For a compact nondegenerate interval \(I\) and \(k\geq0\), put
    \[
        \|u\|_{C^k(I)}=\sum_{j=0}^k\|u^{(j)}\|_{\infty,I},
        \qquad u\in C^k(I,\realnumbers^m).
    \]
\end{snippetdefinition}

\begin{snippetproposition}{diffeq-function-spaces-complete}{Completeness and an anchored norm}
    The spaces \(B(A,\realnumbers^m)\) with the uniform norm and
    \(C^k([a,b],\realnumbers^m)\) with the uniform derivative norm are complete.
    For \(x_0\in[a,b]\), the norm
    \[
        \|u\|_*=\|u(x_0)\|+\|u'\|_\infty
    \]
    is equivalent to the usual norm on \(C^1([a,b],\realnumbers^m)\).
\end{snippetproposition}

\begin{snippetproof}{diffeq-function-spaces-complete-proof}{diffeq-function-spaces-complete}{Completeness and an anchored norm}
    A uniformly Cauchy sequence has pointwise limits in the complete space
    \(\realnumbers^m\); the Cauchy estimate passes to the limit uniformly.
    Uniform limits preserve boundedness and continuity.
    If a sequence is Cauchy in \(C^k\), each derivative has a uniform limit.
    Applying the fundamental theorem of calculus successively shows that these
    limits are the derivatives of the zeroth limit, proving completeness.
    Finally
    \[
        \|u\|_*\leq\|u\|_{C^1},\qquad
        \|u\|_\infty\leq\|u(x_0)\|+(b-a)\|u'\|_\infty,
    \]
    so \(\|u\|_{C^1}\leq\max\{1,1+b-a\}\|u\|_*\).
\end{snippetproof}

\section{The integral formulation}

\begin{snippetdefinition}{diffeq-volterra-equation-definition}{Volterra equation}
    Let \(f\colon\Omega\fromto\realnumbers^m\) be continuous on an open
    \(\Omega\subseteq\realnumbers^{m+1}\), and let \((x_0,y_0)\in\Omega\).
    A solution of the associated \emph{Volterra equation} on an interval \(I\)
    containing \(x_0\) is a continuous \function \(u\colon I\fromto\realnumbers^m\)
    with graph in \(\Omega\), such that
    \[
        u(x)=y_0+\integral[x_0][x][f(t,u(t))][t].
    \]
\end{snippetdefinition}

\begin{snippetproposition}{diffeq-volterra-equivalence}{Differential and integral formulations}
    Let \(f\colon\Omega\fromto\realnumbers^m\) be continuous on an open
    \(\Omega\subseteq\realnumbers^{m+1}\), and let \((x_0,y_0)\in\Omega\).
    A continuous \function \(u\) on an interval \(I\ni x_0\), with graph in
    \(\Omega\), solves the Volterra equation
    \[
        u(x)=y_0+\int_{x_0}^x f(t,u(t))\,dt
    \]
    \ifandonlyif it is a \(C^1\) solution of \(u'=f(x,u)\), \(u(x_0)=y_0\).
\end{snippetproposition}

\begin{snippetproof}{diffeq-volterra-equivalence-proof}{diffeq-volterra-equivalence}{Differential and integral formulations}
    For a differential solution, the fundamental theorem of calculus gives
    \[
        u(x)-y_0=u(x)-u(x_0)=\int_{x_0}^x u'(t)\,dt
        =\int_{x_0}^x f(t,u(t))\,dt.
    \]
    Conversely \(t\mapsto f(t,u(t))\) is continuous, so its proper Riemann
    integral is \(C^1\) with derivative \(f(x,u(x))\). Evaluating at \(x_0\)
    gives the initial condition.
\end{snippetproof}

\begin{snippetproposition}{diffeq-uniform-limit-solutions}{Uniform limits of solutions}
    Let \(f\colon\Omega\fromto\realnumbers^m\) be continuous on an open set,
    and let \(u_n\) solve \(u_n'=f(x,u_n)\) on a compact interval \(I\).
    If \(u_n\to u\) uniformly and all graphs, including that of \(u\),
    lie in one compact subset of \(\Omega\), then \(u\) solves the same equation
    and \(u_n'\to u'\) uniformly.
    In particular the compact-graph condition is automatic when
    \(\Omega=\realnumbers^{m+1}\).
\end{snippetproposition}

\begin{snippetproof}{diffeq-uniform-limit-solutions-proof}{diffeq-uniform-limit-solutions}{Uniform limits of solutions}
    Uniform continuity of \(f\) on the common compact set gives
    \(f(\cdot,u_n)\to f(\cdot,u)\) uniformly.
    In the identities
    \(u_n(x)=u_n(x_0)+\int_{x_0}^x f(t,u_n(t))\,dt\),
    pass to the uniform limit. The Volterra equivalence proves the result.
    Equivalently, uniform convergence of the derivatives and convergence at
    one point allow passage to the limit under differentiation.
\end{snippetproof}
\end{document}
